\documentclass[11pt]{article}
\usepackage{fontspec}
\usepackage{xeCJK}
\usepackage{amsmath}
\usepackage{amssymb}
\usepackage{array}
\usepackage{booktabs}
\usepackage[margin=1in]{geometry}
\usepackage[hidelinks]{hyperref}
\usepackage{seqsplit}

\setmainfont{Times New Roman}
\setCJKmainfont{Malgun Gothic}
\setCJKsansfont{Malgun Gothic}
\setCJKmonofont{Malgun Gothic}
\xeCJKsetup{CJKspace=true}
\XeTeXlinebreaklocale "ko"
\XeTeXlinebreakskip=0pt plus 1pt
\setlength{\emergencystretch}{2em}

\newcommand{\claimstatus}[2]{\textbf{#1 [#2]}}
\newcommand{\poly}{\operatorname{poly}}
\newcommand{\bits}{\operatorname{bitlength}}
\newcommand{\archivepaper}{\path{paper/main-ko.tex}}
\newcolumntype{L}[1]{>{\raggedright\arraybackslash}p{#1}}
\newenvironment{claimproof}{\par\noindent\textbf{증명.}}{\hfill\(\square\)\par}
\newenvironment{costmodel}[1]{%
  \begin{center}
  \small
  \begin{tabular}{L{0.19\linewidth}L{0.69\linewidth}}
  \toprule
  \multicolumn{2}{l}{\textbf{#1}}\\
  \midrule
}{%
  \bottomrule
  \end{tabular}
  \end{center}
}
\newcommand{\costrow}[2]{\textbf{#1} & #2\tabularnewline}
\renewcommand{\abstractname}{초록}


\title{부호 있는 기하급수 회로의\\
예외적 Cyclotomic 추출}
\author{MOSEF 연구 프로그램}
\date{2026년 7월 31일}

\begin{document}
\maketitle

\begin{abstract}
본 논문은 하나의 부호 있는 깊이 2 기하급수 회로에서 나타나는 유리
root-of-unity 비율을 완전히 분류하고 두 개의 예외적 이차 cyclotomic
족을 분리한다. 공통 step 경우를 제외하면 비율 1인 차수 4와 비율 2인
차수 6만 남는다. 두 quotient cofactor는 나눗셈 없이 로그 개수의
modular step으로 평가할 수 있고, 직접 cyclotomic, stage, 공개 overlap
GCD가 모두 unit인 뒤에도 proper factor를 드러낼 수 있다. 정확한
resultant 공식은 모든 overlap branch를 공개적이고 total하게 만든다.
동시에 고정된 유한 schedule은 무한히 많은 semiprime을 놓치며, exact
lift materialization은 compact 평가보다 지수적으로 비쌀 수 있고,
한 cofactor는 signature 한 비트만 준다는 한계를 증명한다. 이는 특정
회로 문법의 결과이며 보편 schedule이나 일반 인수분해 알고리즘이 아니다.
\end{abstract}

\section{비용이 명시된 회로 문법}

\[
S_M(X)=1+X+\cdots+X^{M-1},\quad
Q_1=S_A(X),\quad Q_2=S_B(X^A)
\]
로 두고, 서로 다른 공개 \(A,B\ge2\)와 0이 아닌 공개 계수에 대해
\[
F(X)=c_1Q_1(X)+c_2Q_2(X)
\]
를 정의한다. 구성, 두 stage 평가, 계수 축약, 모든 GCD, 실패한 역원,
독립적인 cyclotomic/cofactor 출구, 요청된 dense/sparse 출력을 모두
비용에 포함한다. 복잡도는
\(\bits(N)=\lfloor\log_2N\rfloor+1\)에 대해 계산한다.

\claimstatus{BAR-018}{PROVED}. 두 stage는 \(\mathbb Q[X]\)에서 서로
소이고
\[
\operatorname{Res}(Q_1,Q_2)=B^{A-1}.
\]
\(L=Q_1(g)\)가 \(N\)의 unit이면
\[
\gcd(F(g),N)=
\gcd(c_1+c_2Q_2(g)L^{-1},N).
\]
proper \(L\)-GCD는 이미 \(N\)을 분해하고, full \(L\)-collision에서는
aggregate가 공개 값 \(c_2B\)로 축약된다. 또한
\(h=\gcd(A-1,B-1)\)에 대해
\[
Q_2(X)-Q_1(X)=XS_h(X)C_{A,B}(X)
\]
이며 두 자연스러운 endpoint와의 다항식 GCD는 정확히 \(S_h(X)\)이다.

\begin{claimproof}
\[
Q_2(X)-B=(X-1)Q_1(X)\sum_{j=1}^{B-1}S_j(X^A)
\]
를 \(A\)-th root에서 평가하면 서로 소성과 resultant가 나온다. 같은
항등식을 \(g\)에서 평가하면 세 modular branch가 완전히 분리된다.
두 endpoint의 root set을 비교하고 monic division을 적용하면 공통
인수 \(S_h\)가 정확히 얻어진다. 따라서 묵시적인 나눗셈은 없다.
\end{claimproof}

\section{Primitive residue와 공개 overlap}

\(d=\gcd(|c_1|,|c_2|)\), \(c_i=da_i\),
\(\gcd(|a_1|,|a_2|)=1\)로 둔다.

\claimstatus{BAR-019}{PROVED}. content \(d\)는 unit, proper, full의
total GCD 출구를 갖는다. unit branch에서 primitive aggregate
\(P=a_1Q_1+a_2Q_2\)의 stage resultant는
\[
\operatorname{Res}(Q_1,F)=(c_2B)^{A-1},\qquad
\operatorname{Res}(Q_2,F)=c_1^{A(B-1)}B^{A-1}
\]
이고,
\[
\gcd(Q_1(g),F(g),N)\mid\gcd(c_2B,N),\quad
\gcd(Q_2(g),F(g),N)\mid\gcd(c_1B,N).
\]
두 stage의 zero set 밖의 root of unity에서는
\[
F(\zeta)=0\iff
c_1(\zeta^A-1)^2+
c_2(\zeta-1)(\zeta^{AB}-1)=0.
\]

\begin{claimproof}
content 분해가 첫 trichotomy를 준다. 각 stage root에서 평가하고
resultant의 곱셈성을 사용하면 두 거듭제곱이 나온다. Bézout 항등식을
\(N\)의 각 소인수로 축약하면 overlap은 공개 coefficient--multiplier
GCD를 나눈다. 유리 stage 식에 두 denominator를 곱하면 마지막 조건이
된다.
\end{claimproof}

\((A,B,c_1,c_2)=(3,7,1,1)\)에서는 \(\Phi_4\mid F\)이다. \(N=55,g=2\)
에서 두 stage와 두 공개 overlap bound가 모두 unit이지만
\(\gcd(F(2),55)=5\)이므로 경계 인수만으로는 충분하지 않다.

\section{유리 root-orbit의 완전 분류}

두 stage의 zero set 밖에 있는 차수 \(n>1\)의 primitive root \(\zeta\)에
대해
\[
R_{A,B}(\zeta)=-\frac{S_B(\zeta^A)}{S_A(\zeta)}
\]
로 둔다.

\claimstatus{THM-003}{PROVED}. 이 비율이 유리수인 경우는 정확히 다음
세 가지이다.
\[
\begin{array}{c|c|c}
\text{차수 조건}&R_{A,B}(\zeta)&(c_1,c_2)\\ \hline
n\mid\gcd(A-1,B-1)&-1&(-1,1)\\
n=4,\ A\equiv B\equiv3\pmod4&1&(1,1)\\
n=6,\ A\equiv5,\ B\equiv3\pmod6&2&(2,1).
\end{array}
\]
요청된 차수의 인식은 \(A,B,n\)의 이진 길이에 대해 다항시간이다.
공통 step 족은 \(\gcd(A-1,B-1)\)의 약수라는 compact descriptor로
남기며, 실제 약수 열거 비용은 별도로 계산한다.

\begin{claimproof}
Galois 불변성과 켤레는 먼저 \(n\mid A(B-2)+1\)을 강제한다.
\(bA\equiv1\pmod n\), \(k\equiv b-1\pmod n\)를 택하면
\[
R_{A,B}(\zeta)+1=
\frac{\sin^2(k\pi/n)}{\sin^2(\pi/n)}
\]
이다. 이는 유리 대수적 정수이므로 음이 아닌 정수이다. cyclotomic norm
case split은 공통 step orbit와 두 이차 orbit만 남긴다. 직접 대입으로
비율과 합동 조건을 확인한다. 위상 합동만으로는 부족하며
\((A,B,n)=(2,4,5)\)에서는 \((1+\sqrt5)/2\)가 나온다. 전체 증명은
\path{research/proofs/THM-003-rational-root-orbits.md}에 있다.
\end{claimproof}

\section{나눗셈 없는 예외적 cofactor}

\[
\begin{aligned}
F_4&=S_A(X)+S_B(X^A)=\Phi_4(X)C_4(X),
&&A\equiv B\equiv3\pmod4,\\
F_6&=2S_A(X)+S_B(X^A)=\Phi_6(X)C_6(X),
&&A\equiv5,\ B\equiv3\pmod6
\end{aligned}
\]
로 정의한다.

\claimstatus{BAR-020}{PROVED}. 각 \(C_i\)는 차수 \(A(B-1)-2\)인 monic
정수 다항식이지만, \(C_i(g)\bmod N\)은 상수 크기 descriptor와
\(O(\log A+\log B)\)회의 modular 연산으로 나눗셈 없이 평가된다.

\begin{claimproof}
\(A=4a+3\), \(B=4b+3\)에 대해
\[
U_4(4t+3;x)=(1+x)S_t(x^4)+x^{4t},\quad
D_4(A;x)=S_A(-x^2)
\]
로 두고 \(r=A-2\), \(s=A(B-2)\), \(k=(s-r)/2\)라 하면
\[
C_4=U_4(A;x)+D_4(A;x)U_4(B;x^A)+x^rS_k(-x^2).
\]
차수 6에서는 \(x^2-x+1\)로 나눈 impulse response가 주기 6이므로
\(D_6(A;x)=\Phi_6(x^A)/\Phi_6(x)\)는 두 개의 주기 6 구간을 갖는다.
\(H=x^3+2x^2+2x+1\), \(J=x^4+x^3-x-1\)라 하면
\[
\begin{split}
C_6={}&2H(x)S_{a+1}(x^6)\\
&+D_6(A;x)(H(x^A)S_b(x^{6A})+x^{6Ab})\\
&+2x^AJ(x)S_{Ab}(x^6).
\end{split}
\]
이진 기하급수 평가가 비용을 증명한다. 정수 곱 항등식은
prime-power valuation의 가법성을 주므로 direct와 cofactor GCD를
독립적으로 사용할 수 있다. dense 출력은 여전히
\(A(B-1)-1\)개 계수 비용을 낸다.
\end{claimproof}

\section{Overlap 및 schedule 장벽}

\claimstatus{BAR-021}{PROVED}. cofactor GCD는 capped 지역 valuation
profile이 0이 아니고 full도 아닐 때와 그때에만 proper이다.
stage/cofactor resultant의 소수 support는 차수 4에서 \(B\), 차수 6에서
\(2B\)에 포함된다. direct/cofactor overlap은
\[
\begin{aligned}
(u_4,v_4)&=\left(\frac{A(B+2)+1}{4},
                 \frac{A(B-2)+1}{4}\right),&
R_4&=u_4^2+v_4^2,\\
(u_6,v_6)&=\left(-\frac{2(A(B-2)+1)}3,
                 \frac{A(B+4)+4}3\right),&
R_6&=u_6^2+u_6v_6+v_6^2
\end{aligned}
\]
로 공개된다. 모든 고정된 유한 public joint schedule은 무한히 많은
square-free semiprime을 놓친다.

\begin{claimproof}
유일분해가 valuation 조건을 준다. 제거된 이차식으로 cofactor를
축약하면 \(u_i+v_iX\)가 되고 그 이차 norm이 \(R_i>0\)이다. 고정
schedule의 모든 양의 charged exact 값을 곱하면 유한 개 소수만 나눈다.
그 support 밖의 서로 다른 두 소수를 택하면 모든 charged GCD가 1이다.
schedule을 입력보다 먼저 고정하는 양화 순서가 핵심이다.
\end{claimproof}

\claimstatus{BAR-022}{PROVED}. balanced population
\[
\mathcal P_m=\{p\text{ prime}:2^{m-1}\le p^2<2^m\}
\]
에서 materialized exact lift의 총 비트 예산이 \(W_m\), hit 소수 수가
\(h_m\), 전체 소수 수가 \(s_m\)이면 정확히
\(\binom{s_m-h_m}{2}\)개의 쌍에서 모든 materialized GCD가 1이고
\[
\left\lfloor\frac{m-1}{2}\right\rfloor h_m\le W_m.
\]
모든 쌍을 다루려면 \(h_m\ge s_m-1\)이어야 한다.

\begin{claimproof}
두 unhit 소수의 곱은 어떤 lift도 나누지 않으므로 모든 GCD가 1이다.
hit 소수의 square-free 곱은 모든 nonzero lift의 곱을 나누고 각 소수는
적어도 \(\lfloor(m-1)/2\rfloor\)비트이므로 예산 부등식이 나온다. 이
materialization 정리는 compact 평가에 그대로 이전되지 않는다.
\end{claimproof}

\claimstatus{BAR-023}{PROVED}. compact 차수 4 tuple
\(A=3\), \(B_m=2^m+3\), \(g=2\)에서 \(C_m=C_4(2)\),
\(E_m=3\cdot2^m+5\)라 하면
\[
C_m=\frac{16(2^{E_m}+3)}{35},\qquad
\gcd(C_m,C_{m+1})=16.
\]
\(s\)개 홀수 소수 중 정확히 \(h\)개가 \(C_m\)을 나누면 한 GCD는
\(h(s-h)\)개 proper 결과, \(\binom h2\)개 full collision,
\(\binom{s-h}{2}\)개 unit 결과를 갖는다. \(s\ge3\)이면 모든 쌍을
다룰 수 없다.

\begin{claimproof}
정확한 \(C_4\) 항등식에 대입하면 closed form이 나온다. 연속 numerator의
Euclidean 축약과 5, 7의 별도 처리가 GCD 16을 준다. 두 support bit가
다를 때와 그때에만 쌍이 분리되므로 세 개수가 따르고
\(h(s-h)<\binom s2\)이다.
\end{claimproof}

\section{1차 문헌 포지셔닝}

\paragraph{M83-R04: 곱셈 schedule.}
여러 지정 거듭제곱을 계산하는 공유 곱셈 schedule은 고전적
addition-chain 문맥이다 \cite{yao1976evaluation}. 이 문헌은 공식
초록만 확인했으므로 증명, 하한, 계산 모형, 상수를 가져오는 데 쓰지
않는다. BAR-018과 BAR-019의 지수 성장, 분모 분기, exact-output
과금은 자기완결 증명을 유지한다.

\paragraph{M83-R05: 근의 단위원과 유리 관계.}
Conway--Jones는 Galois automorphism, 짧은 root-of-unity 소멸합,
유리 cosine 관계를 확립된 도구로 사용한다
\cite{conway1976trigonometric}. THM-003과 BAR-020은 제한된 signed
depth-two 관측량에 대해 common-step/\(\Phi_4\)/\(\Phi_6\) 삼분류와
나눗셈 없는 evaluator를 재구성한다. 이는 확립된 방법족 안의 제한적
포지셔닝이며 독창성이나 우선권 주장이 아니다.

등록된 M83 행렬은 전체 문헌 경계를 보존한다. 이 제한 감사에서 같은
식을 찾지 못한 사실은 독창성의 증거가 아니다.

\section{재현성과 한계}

\begin{verbatim}
python scripts/run_m24_rational_residue_audit.py
python scripts/check_m24_rational_residue_audit_differential.py
python scripts/run_m25_rational_root_orbit_audit.py
python scripts/check_m25_rational_root_orbit_differential.py
python scripts/run_m26_exceptional_cyclotomic_audit.py
python scripts/check_m26_exceptional_cyclotomic_differential.py
\end{verbatim}
EXP-0023부터 EXP-0025까지의 기록은
\path{research/experiments}에 있다. 유한 symbolic/modular 검사는
구현을 검증하며 대수적 증명을 대체하지 않는다.

모든 전체 증명과 반례 및 branch 비용은 \archivepaper{}에 보존된다.
claim 상태와 집중 논문 투영은 \path{research/CLAIMS.md}와
\path{schemas/m82-paper-portfolio-v1.json}에 있다. 본 논문은 표시된
부호 있는 기하급수 문법만 다루며, universal selector, prime-support
density 정리, 일반 인수분해 알고리즘을 제공하지 않는다.

\bibliographystyle{plain}
\small
\bibliography{paper/references}

\end{document}
